% ============================================================
% ПРЕАМБУЛА
% ============================================================
% Преамбула — это служебная часть документа, которая настраивает
% его внешний вид и подключает необходимые пакеты. Здесь мы:
%   1. Задаём кодировку и язык документа.
%   2. Устанавливаем геометрию страницы (поля, размер бумаги).
%   3. Подключаем пакеты для работы с математикой, графикой,
%      листингами кода, гиперссылками и цветами.
%   4. Настраиваем стиль отображения кода Python.
%   5. Определяем команды для оформления примечаний и предупреждений.


% Кодировка и язык
\usepackage[utf8]{inputenc}           % поддержка UTF-8 (русские буквы)
\usepackage[T2A]{fontenc}
\usepackage[russian]{babel}           % русский язык в документе

% Математика и символы
\usepackage{amsmath, amssymb}         % расширенные математические возможности
\usepackage{bm}                       % жирные символы в математике

% Геометрия страницы
\usepackage{geometry}
\geometry{
	top=2cm,
	bottom=2cm,
	left=2.5cm,
	right=2.5cm
}

% Графика и плавающие объекты
\usepackage{graphicx}                 % вставка изображений
\usepackage{float}                    % управление плавающими объектами
\usepackage{caption}                  % настройка подписей
\usepackage{subcaption}               % подписи для подрисунков

% Списки
\usepackage{enumitem}                 % гибкие списки

% Цвета
\usepackage{color}
\usepackage{xcolor}                   % расширенная работа с цветами

% Гиперссылки
\usepackage{hyperref}
\hypersetup{
	colorlinks=true,
	linkcolor=blue,
	urlcolor=blue,
	citecolor=blue
}


\usepackage[dvipsnames]{xcolor}

% Команды для оформления примечаний
\newcommand{\note}[1]{\textcolor{RoyalBlue}{\textbf{Примечание:} #1}}
\newcommand{\warning}[1]{\textcolor{BurntOrange}{\textbf{Важно:} #1}}
\newcommand{\tip}[1]{\textcolor{ForestGreen}{\textbf{Совет:} #1}}
